\newpage
% ======================= PARTE E =======================
\section{Parte E — Estilo del examen real (2025)}

\begin{nota}
\textbf{Qué muestra el examen de 2025.} Cuatro preguntas de 25 puntos: (1) modelar un ambiente real y
reformularlo como juego competitivo; (2) A* en una gráfica pequeña con heurística admisible pero
inconsistente; (3) modelar un juego \emph{estocástico} y definir su operador de valor
(expectiminimax); (4) leer una KB en CNF de primer orden, derivar por resolución en 2–3 pasos y decidir
si otra conclusión se sigue. Se permitía \textbf{una hoja A4 de notas}. Predominan el modelado y la
explicación sobre las demostraciones largas: responde con estructura (listas, tablas) y justifica cada
decisión en una o dos líneas.
\end{nota}

\begin{ejercicio}{E1 · Examen 1 de 2025 (resuélvelo en 2 h) \dif{2}\fuente{Examen del profesor, 2025}}
\textbf{1. Agentes y ambientes competitivos (25).} Te contratan para diseñar un sistema de taxis
autónomos en la Ciudad de México; la tarea es decidir dónde colocar los taxis.
(a) \pts{15} Modela el ambiente: estados, perceptos, acciones y propiedades (observable/parcial,
determinista/estocástico, discreto/continuo). Sé breve.
(b) \pts{10} Existen compañías que compiten por el espacio. Reformula como juego competitivo y discute
qué algoritmo (minimax, expectimax) sería más adecuado.

\textbf{2. Búsqueda informada (25).} Aristas $S\to A\,(1)$, $S\to B\,(4)$, $A\to C\,(2)$, $B\to C\,(1)$, $C\to G\,(3)$;
$h(S)=5$, $h(A)=3$, $h(B)=2$, $h(C)=1$, $h(G)=0$.
(a) \pts{10} Ejecuta A* mostrando el orden de expansión. (b) \pts{10} ¿La solución es óptima? Justifica.
(c) \pts{5} La heurística es admisible pero no consistente. Explica por qué y qué cambios de
implementación preservan la optimalidad.

\textbf{3. Juegos con incertidumbre (25).} Connect-4 de $4\times4$: al dejar caer una ficha en una columna,
con probabilidad $\rho$ cae en esa columna y con probabilidad $1-\rho$ en una columna uniformemente aleatoria.
(a) \pts{10} Modela formalmente: estados, acciones, transición estocástica.
(b) \pts{15} Define un operador de valor que combine minimax y expectimax para este juego.

\textbf{4. Lógica: CNF y resolución (25).} KB:
1.\ $\neg\p{Estudiante}(x)\vee\neg\p{TomaIA}(x)\vee\p{Conoce}(x,\p{astar})$;\;
2.\ $\p{Algoritmo}(\p{astar})$;\; 3.\ $\p{Optimo}(\p{astar})$;\;
4.\ $\neg\p{Optimo}(a)\vee\p{UsaAdmisible}(a)$;\; 5.\ $\neg\p{UsaAdmisible}(a)\vee\p{Correcto}(a)$;\;
6.\ $\p{Estudiante}(\p{bob})$;\; 7.\ $\p{TomaIA}(\p{bob})$.
(a) \pts{10} Lectura en lenguaje natural de cada cláusula. (b) \pts{10} Por resolución, deriva
$\p{Conoce}(\p{bob},\p{astar})$ en 2–3 pasos. (c) \pts{5} ¿Se sigue $\p{Correcto}(\p{astar})$? Si sí, argumento
breve por resolución; si no, qué falta en la base.
\end{ejercicio}

\begin{ejercicio}{E2 · Modelar un ambiente real y su versión competitiva \dif{2}\fuente{estilo 2025, pregunta 1}}
Diseñas el agente que redistribuye bicicletas entre las estaciones de un sistema de bicis compartidas
(camiones que mueven bicis de estaciones llenas a vacías).
(a) Modela el ambiente: medida de desempeño, estados, perceptos, acciones y las seis propiedades
(observable, agentes, determinista, episódico, estático, discreto). ¿Es un problema de búsqueda, de
CSP o de búsqueda local? Justifica.
(b) Una empresa rival de scooters compite por los mismos usuarios y también decide dónde colocar sus
vehículos. Formúlalo como juego (jugadores, turnos, acciones, utilidades). ¿Es de suma cero? ¿Usarías
minimax, expectimax o expectiminimax? ¿Qué harías si el juego es demasiado grande para llegar a las hojas?
\end{ejercicio}

\begin{ejercicio}{E3 · A* con heurística inconsistente y reapertura \dif{2}\fuente{estilo 2025, pregunta 2}}
Aristas dirigidas: $S\to A\,(2)$, $S\to B\,(1)$, $A\to C\,(1)$, $B\to C\,(3)$, $C\to D\,(1)$, $D\to G\,(2)$.
Heurística: $h(S)=5$, $h(A)=4$, $h(B)=1$, $h(C)=1$, $h(D)=2$, $h(G)=0$.
(a) Calcula $h^*$; muestra que $h$ es admisible y encuentra los arcos que violan la consistencia.
(b) Ejecuta A* en grafo \emph{sin} reapertura. ¿Qué costo devuelve? ¿Es óptimo?
(c) Ejecuta A* \emph{con} reapertura de nodos cerrados. ¿Qué cambia?
(d) ¿Cuál es el mínimo número de valores de $h$ que hay que cambiar para que sea consistente (y siga admisible)? Propón uno.
\end{ejercicio}

\begin{ejercicio}{E4 · Gato con viento \dif{2}\fuente{estilo 2025, pregunta 3}}
En un gato estocástico, al marcar una casilla libre, con probabilidad $\rho$ la marca queda ahí y con
probabilidad $1-\rho$ queda en una casilla libre distinta, elegida uniformemente.
(a) Modela formalmente: estados, jugador en turno, acciones, modelo de transición $P(s'\mid s,a)$,
prueba terminal y utilidad.
(b) Escribe el operador expectiminimax para este juego con profundidad limitada y función de evaluación.
(c) En cierta posición MAX tiene dos jugadas. \emph{Atacar}: si la marca cae donde quiere, gana ($+1$);
si se desvía, pierde ($-1$). \emph{Bloquear}: si cae donde quiere, empata ($0$); si se desvía, cae con igual
probabilidad en una casilla que empata ($0$) o en una que pierde ($-1$). Calcula el valor de ambas jugadas
en función de $\rho$ y el umbral de $\rho$ a partir del cual conviene atacar. ¿Qué elegiría minimax (tratando
el azar como adversario)?
\end{ejercicio}

\begin{ejercicio}{E5 · KB en CNF: leer, derivar y detectar qué falta \dif{2}\fuente{estilo 2025, pregunta 4}}
KB: 1.\ $\neg\p{Robot}(x)\vee\neg\p{Bateria}(x)\vee\p{Mueve}(x)$;\; 2.\ $\neg\p{Mueve}(x)\vee\neg\p{EnRuta}(x)\vee\p{Entrega}(x)$;\;
3.\ $\p{Robot}(r_1)$;\; 4.\ $\p{Bateria}(r_1)$;\; 5.\ $\neg\p{Dron}(x)\vee\p{Vuela}(x)$;\; 6.\ $\neg\p{Vuela}(x)\vee\p{EvitaTrafico}(x)$;\;
7.\ $\p{EnRuta}(r_2)$.
(a) Lee cada cláusula en lenguaje natural (como implicación).
(b) Deriva $\p{Mueve}(r_1)$ por resolución en 2 pasos.
(c) ¿Se sigue $\p{Entrega}(r_1)$? ¿Y $\p{EvitaTrafico}(r_1)$? Si no, di qué falta y da un modelo de la KB donde sea falsa.
(d) Pasa a CNF: «Todo dron con batería baja regresa a la base» y «Todo pedido es entregado por algún
robot». ¿Cuál necesita skolemización?
\end{ejercicio}
